% Half-title, title and copyright pages.

\thispagestyle{empty}
\vspace*{2.2in}
\begin{flushright}
  {\Large\bfseries Modern Embedded C++}
\end{flushright}
\cleardoublepage

\thispagestyle{empty}
\vspace*{1.1in}
\begin{flushleft}
  {\color{accent}\rule{0.9in}{2.5pt}}\par\vspace{1.4ex}
  {\fontsize{32}{38}\selectfont\bfseries Modern\\[0.15ex] Embedded C++\par}
  \vspace{2.2ex}
  {\Large\itshape From C structs to classes, interfaces,\\templates and threads\par}
  \vspace{1.2in}
  {\large Erik Pihl\par}
  \vfill
  {\small\scshape\lsstyle qrtech academy\par}
  \vspace{0.6ex}
  {\small\color{muted} Draft of 10 September 2026\par}
\end{flushleft}
\newpage

\thispagestyle{empty}
\vspace*{\fill}
{\small
\noindent\textit{Modern Embedded C++}\\
Copyright \textcopyright\ 2026 QRTECH Academy.

\medskip
\noindent The course this book is typeset from, and every line of code in it, is released under
the MIT License. The companion repository holds the lecture material, the example programs and
the solutions to every exercise, which Appendix~\ref{sol:ch:solutions} typesets straight from
their files:

\medskip
\noindent\url{https://github.com/qrtech-academy/modern-embedded-cpp}

\medskip
\noindent The C++ in this book is C++17 throughout, compiled with GCC. ESP32-S3, STM32 and
ATmega328P appear in the text as examples of targets a driver might be written for; they are
trademarks of their respective owners.

\medskip
\noindent Set in TeX Gyre Pagella and DejaVu Sans Mono with Lua\LaTeX.

\medskip
\noindent\textcolor{muted}{This is a draft, typeset from the six lectures of the course, their
exercises and solutions, and the two self-assessment papers.}
\par}
\cleardoublepage
